\documentclass[10pt,a4paper]{amsart}
\usepackage[margin=19mm]{geometry}
\usepackage{amsmath,amssymb,mathtools}
\usepackage[scr=rsfs,cal=euler]{mathalfa}
\usepackage{xfrac}
\usepackage[unicode,hidelinks,bookmarks=false]{hyperref}
\newcommand{\ie}{i.e.\ }
\newcommand{\cf}{cf.\ }
\newcommand{\resp}{respectively\ }
\newcommand{\Subsection}{\subsection}
\providecommand{\nobreakdash}{\nobreak\hbox{-}}
\setlength{\emergencystretch}{2em}
\multlinegap0pt
\usepackage{amssymb}
\usepackage{xcolor}
\usepackage{bm}
\newtheorem{lemma}{Lemma}
\newtheorem{theorem}{Theorem}
\renewcommand{\ge}{\geqslant}
\renewcommand{\geq}{\geqslant}
\renewcommand{\hat}{\widehat}
\renewcommand{\leq}{\leqslant}
\title{Gradient-free stochastic optimization of derivatives under strong convexity}
\author{Arya Akhavan, Sirine Louati, Alexandre B. Tsybakov}
\date{Source: arXiv:2607.07249v1 (CC BY 4.0)}
\begin{document}
\maketitle

\noindent\emph{Source and licence.} Gradient-free stochastic optimization of derivatives under strong convexity. Version arXiv:2607.07249v1 (CC BY 4.0), distributed by arXiv under the Creative Commons Attribution 4.0 International licence, http://creativecommons.org/licenses/by/4.0/, as recorded in the arXiv licence metadata for this exact version and shown on the abstract page https://arxiv.org/abs/2607.07249v1. Full licence notice: \texttt{licenses/arxiv-2607.07249-CC-BY-4.0.txt}.\par\smallskip

\noindent\textit{Editorial scope.} Counted unit: the article's theorem thm:upper-bound (source0.tex lines 800--819). The article's complete original proof of it is delivered with it (source0.tex lines 1317--1538, one \texttt{proof} environment of 222 lines, which the article places away from the statement). No mathematics is written, repaired or re-derived: the statement and the proof are the article's own lines, in the article's own notation, and no retained line is substituted or re-flowed.  Retained uncounted as the source-local context the proof needs: lines 641--645; lines 655--659; lines 749--768; lines 770--785; lines 1201--1217.  Nothing was omitted from any retained span, so the package carries no omission record.  No retained line contains a citation, so the wrapper carries no bibliography and no bibliography entry.  No retained line contains a figure environment, an image-inclusion command or drawing code, so no image is delivered and the package declares no asset dependency.  Editorial formatting only, in addition: an AMS \texttt{amsart} preamble that restates the article's own declarations and macros used by the retained lines, namely the environments \texttt{lemma}, \texttt{theorem} on separate counters, together with the standard packages those lines need.  The retained environments print in this document as Lemma 1, Theorem 1 in order of appearance, whereas the article prints the same items with the numbering of its own full document.  The complete native source of the article as distributed in the arXiv e-print for this exact version is bundled unchanged as \texttt{sources/204691/source0.tex}.
\begin{equation}
\label{eq:smoothness}
f(x)-f(x^\star)\leq\frac{L_f}{2}\|x-x^\star\|^2,
\qquad\forall\,x\in\Theta.
\end{equation}

\begin{equation}
\label{eq:sgd}
x_{t+1}:=\Pi_\Theta\bigl(x_t-\eta_t\,\tilde{g}_t\bigr),
\qquad t=0,1,\dots,T-1,
\end{equation}

\begin{lemma}[Bias of the gradient estimator]
\label{lem:bias-multi}
Let $g\in\mathcal{F}_{\beta+k}(L)$ with $\beta>1$ and $k\geq 0$. 
For every $x_t\in\mathbb{R}^d$ and every $l\in\{1,\dots,d\}$,
\[
\bigl|\mathbb{E}[\tilde{g}_{t,l}\mid x_t]-\partial_l f(x_t)\bigr|
\leq C_{\mathrm{bias}}\,h^{\beta-1},
\]
where 
\[
C_{\mathrm{bias}}
:=L\,\bigl(\mathbb{E}[K_k(U)^2]\bigr)^{1/2}
\bigl(\mathbb{E}[K_1(V)^2]\bigr)^{1/2}.
\]
Consequently,
\[
\bigl\|\mathbb{E}[\tilde{g}_t\mid x_t]-\nabla f(x_t)\bigr\|^2
\leq d\,C_{\mathrm{bias}}^2\,h^{2(\beta-1)}.
\]
\end{lemma}

\begin{lemma}[Second moment of the gradient estimator]
\label{lem:var-multi-second-moment}
Let $g\in\mathcal{F}'_{\alpha,\beta,k}(L,L_f,G,G_g,\bar{L})$. For 
any $x_t\in\Theta$ and $h\in(0,1]$,
\[
\mathbb{E}\bigl[\|\tilde{g}_t\|^2\mid x_t\bigr]
\leq 4G^2+4d\,C_{\mathrm{bias}}^2\,h^{2(\beta-1)}
+2d\,C_{\mathrm{var}}\,h^{-2(k+1)},
\]
where 
\[
C_{\mathrm{var}}
:=\frac{\mathbb{E}[K_k(U)^2]\,\mathbb{E}[K_1(V)^2]}{4}
\Bigl(\frac{3\bar{L}^2}{16}+6G_g^2+8\sigma^2\Bigr).
\]
\end{lemma}

\begin{lemma}[Recursive inequality]
\label{lem:recursion}
Let $(u_t)_{t\geq 0}$ be a sequence of non-negative real numbers 
satisfying
\begin{equation}
\label{eq:recursion-hypo}
u_{t+1}\leq\Bigl(1-\frac{q}{t+t_0}\Bigr)u_t
+\frac{A}{t+t_0}+\frac{B}{(t+t_0)^2},\qquad t\geq 0,
\end{equation}
for some constants $q\geq 2$, $t_0\geq 2q$ and $A,B\geq 0$. Then, 
for all $T\geq 1$,
\begin{equation}
\label{eq:recursion-conclusion}
u_T\leq\frac{t_0\,u_0}{T+t_0}+\frac{2A}{q}
+\frac{2B}{(q-1)(T+t_0)}.
\end{equation}
\end{lemma}

\begin{theorem}[Upper bound]
\label{thm:upper-bound}
Let $\beta\ge 2$, $k\geq 0$, and let $\alpha,L,L_f,G,G_g,\bar{L}>0$ 
be given constants. Consider the projected stochastic gradient 
algorithm~\eqref{eq:sgd} with step size $\eta_t=\gamma/(t+t_0)$, 
where $\gamma\geq 4/\alpha$ and $t_0\geq\alpha\gamma$ and smoothing
parameter $h=\min\bigl(\kappa\,T^{-1/(2(\beta+k))},1\bigr)$ for 
some $\kappa>0$. Set $N=2dT$. Then, for any initialization 
$x_0\in\Theta$,
\[
\sup_{g\in\mathcal{F}'_{\alpha,\beta,k}(L,L_f,G,G_g,\bar{L})}
\mathbb{E}\bigl[f(x_T)-f(x^\star)\bigr]
\leq C\,d^{(2\beta+k-1)/(\beta+k)}\,N^{-(\beta-1)/(\beta+k)},
\]
where $C>0$ is a constant depending only on 
$\alpha,\gamma,t_0,\kappa,L_f,G,R,C_{\mathrm{bias}},C_{\mathrm{var}}$ 
and independent of $N$ and $d$. %In particular, the iterate 
%$x_T=\hat{x}_N$ produced by the algorithm after $T=N/(2d)$ steps 
%attains the rate above.
\end{theorem}

\begin{proof}[Proof of Theorem~\ref{thm:upper-bound}]
\label{proof:thm1}
Throughout the proof, fix 
$g\in\mathcal{F}'_{\alpha,\beta,k}(L,L_f,G,G_g,\bar{L})$ and 
let $x^\star\in\mathrm{int}(\Theta)$ denote the minimizer of 
$f=\partial_j^k g$ on $\Theta$, so that $\nabla f(x^\star)=0$. The 
recursion~\eqref{eq:sgd} guarantees $x_t\in\Theta$ for all 
$t\geq 0$, hence properties (b), (c), (d) of the class apply at 
every iteration. Set
\[
\Delta_t:=\mathbb{E}\bigl[\|x_t-x^\star\|^2\bigr],
\qquad 
b_t:=\mathbb{E}[\tilde{g}_t\mid x_t]-\nabla f(x_t),
\]
and note that $\Delta_0\leq R^2$.

\medskip
\noindent\textit{One-step recursive inequality on $\Delta_t$.}
Since $\Pi_\Theta$ is non-expansive and $x^\star\in\Theta$,
\[
\|x_{t+1}-x^\star\|^2
\leq\|x_t-\eta_t\tilde{g}_t-x^\star\|^2
=\|x_t-x^\star\|^2-2\eta_t\langle\tilde{g}_t,x_t-x^\star\rangle
+\eta_t^2\|\tilde{g}_t\|^2.
\]
Taking the conditional expectation given $x_t$, using 
$\mathbb{E}[\tilde{g}_t\mid x_t]=\nabla f(x_t)+b_t$, then taking 
total expectation,
\begin{equation}
\label{eq:upper-rec1}
\Delta_{t+1}\leq\Delta_t
-2\eta_t\,\mathbb{E}\bigl[\langle\nabla f(x_t),x_t-x^\star\rangle\bigr]
-2\eta_t\,\mathbb{E}\bigl[\langle b_t,x_t-x^\star\rangle\bigr]
+\eta_t^2\,\mathbb{E}\bigl[\|\tilde{g}_t\|^2\bigr].
\end{equation}
By $\alpha$-strong convexity of $f$ on $\Theta$ and 
$\nabla f(x^\star)=0$,
\begin{equation}
\label{eq:upper-sc}
\langle\nabla f(x_t),x_t-x^\star\rangle
\geq f(x_t)-f(x^\star)+\frac{\alpha}{2}\|x_t-x^\star\|^2
\geq\frac{\alpha}{2}\|x_t-x^\star\|^2.
\end{equation}
By the Cauchy--Schwarz inequality and Young's inequality 
$2|\langle u,v\rangle|\leq\varepsilon\|u\|^2+\|v\|^2/\varepsilon$ 
applied with $\varepsilon=\alpha/2$,
\begin{equation}
\label{eq:upper-bias}
2\bigl|\langle b_t,x_t-x^\star\rangle\bigr|
\leq\frac{\alpha}{2}\|x_t-x^\star\|^2+\frac{2}{\alpha}\|b_t\|^2.
\end{equation}
Substituting~\eqref{eq:upper-sc} and~\eqref{eq:upper-bias} 
into~\eqref{eq:upper-rec1},
\begin{equation}
\label{eq:upper-rec2}
\Delta_{t+1}\leq\Bigl(1-\frac{\alpha\eta_t}{2}\Bigr)\Delta_t
+\frac{2\eta_t}{\alpha}\,\mathbb{E}\bigl[\|b_t\|^2\bigr]
+\eta_t^2\,\mathbb{E}\bigl[\|\tilde{g}_t\|^2\bigr].
\end{equation}

\medskip
\noindent\textit{Reduction to a Chung-type recursion.}
By Lemma~\ref{lem:bias-multi}, 
\begin{equation}
\label{eq:upper-bias-bound}
\mathbb{E}\bigl[\|b_t\|^2\bigr]
\leq d\,C_{\mathrm{bias}}^2\,h^{2(\beta-1)},
\end{equation}
and by Lemma~\ref{lem:var-multi-second-moment}, applicable since 
$x_t\in\Theta$ and $h\in(0,1]$,
\begin{equation}
\label{eq:upper-var-bound}
\mathbb{E}\bigl[\|\tilde{g}_t\|^2\bigr]
\leq 4G^2+4d\,C_{\mathrm{bias}}^2\,h^{2(\beta-1)}
+2d\,C_{\mathrm{var}}\,h^{-2(k+1)}.
\end{equation}
Substituting~\eqref{eq:upper-bias-bound} 
and~\eqref{eq:upper-var-bound} into~\eqref{eq:upper-rec2} with 
$\eta_t=\gamma/(t+t_0)$ yields
\begin{equation}
\label{eq:upper-rec3}
\Delta_{t+1}\leq\Bigl(1-\frac{q}{t+t_0}\Bigr)\Delta_t
+\frac{A}{t+t_0}+\frac{B}{(t+t_0)^2},
\end{equation}
where 
\begin{equation}
\label{eq:upper-qAB}
q:=\frac{\alpha\gamma}{2},\quad
A:=\frac{2\gamma\,d\,C_{\mathrm{bias}}^2\,h^{2(\beta-1)}}{\alpha},\quad
B:=\gamma^2\bigl(4G^2+4d\,C_{\mathrm{bias}}^2\,h^{2(\beta-1)}
+2d\,C_{\mathrm{var}}\,h^{-2(k+1)}\bigr).
\end{equation}
The conditions $\gamma\geq 4/\alpha$ and $t_0\geq\alpha\gamma$ 
imposed in the statement of the theorem are equivalent to $q\geq 2$ 
and $t_0\geq 2q$, which are precisely the hypotheses of 
Lemma~\ref{lem:recursion}. Applying this lemma 
to~\eqref{eq:upper-rec3} with $u_t=\Delta_t$ and $u_0\leq R^2$ 
yields, for all $T\geq 1$,
\begin{equation}
\label{eq:upper-DeltaT-from-lemma}
\Delta_T\leq\frac{t_0\,R^2}{T+t_0}
+\frac{2A}{q}
+\frac{2B}{(q-1)(T+t_0)}.
\end{equation}

\medskip
\noindent\textit{Bias-variance bound on $\Delta_T$.}
Using $q=\alpha\gamma/2$ in $2A/q$, and $T+t_0\geq T$ in the third 
term of~\eqref{eq:upper-DeltaT-from-lemma},
\[
\frac{2A}{q}=\frac{8\,d\,C_{\mathrm{bias}}^2\,h^{2(\beta-1)}}{\alpha^2},
\qquad
\frac{2B}{(q-1)(T+t_0)}\leq\frac{2B}{(q-1)T}.
\]
Substituting the expression of $B$ from~\eqref{eq:upper-qAB} and 
grouping terms,
\begin{equation}
\label{eq:upper-DeltaT-grouped}
\Delta_T\leq\frac{C_3}{T}+\frac{C_4\,d\,h^{-2(k+1)}}{T}
+C_5\,d\,h^{2(\beta-1)},
\end{equation}
where 
\begin{equation}
\label{eq:upper-C345}
C_3:=t_0\,R^2+\frac{8\gamma^2 G^2}{q-1},\quad
C_4:=\frac{4\gamma^2\,C_{\mathrm{var}}}{q-1},\quad
C_5:=\frac{8\,C_{\mathrm{bias}}^2}{\alpha^2}
+\frac{8\gamma^2\,C_{\mathrm{bias}}^2}{q-1}.
\end{equation}
The contribution 
$8\gamma^2\,d\,C_{\mathrm{bias}}^2\,h^{2(\beta-1)}/((q-1)T)$ 
arising from the bias-related part of 
$2B/((q-1)T)$ has been absorbed into the term 
$C_5\,d\,h^{2(\beta-1)}$, using $1/T\leq 1$ for $T\geq 1$.

\medskip
\noindent\textit{Conversion to optimization error.}
By the smoothness inequality~\eqref{eq:smoothness} (valid since 
$x_T\in\Theta$ and $\nabla f(x^\star)=0$),
\[
\mathbb{E}\bigl[f(x_T)-f(x^\star)\bigr]
\leq\frac{L_f}{2}\,\Delta_T,
\]
which combined with~\eqref{eq:upper-DeltaT-grouped} gives
\begin{equation}
\label{eq:upper-opt-error}
\mathbb{E}\bigl[f(x_T)-f(x^\star)\bigr]
\leq\frac{C_6}{T}+\frac{C_7\,d\,h^{-2(k+1)}}{T}
+C_8\,d\,h^{2(\beta-1)},
\end{equation}
where $C_i:=L_f\,C_{i-3}/2$ for $i\in\{6,7,8\}$.

\medskip
\noindent\textit{Optimal choice of $h$.}
We now substitute the prescribed value 
$h=\min(\kappa\,T^{-1/(2(\beta+k))},1)$ and consider two regimes.

If $T\geq\kappa^{2(\beta+k)}$, then 
$h=\kappa\,T^{-1/(2(\beta+k))}\leq 1$, so
\[
h^{-2(k+1)}=\kappa^{-2(k+1)}\,T^{(k+1)/(\beta+k)},
\qquad
h^{2(\beta-1)}=\kappa^{2(\beta-1)}\,T^{-(\beta-1)/(\beta+k)}.
\]
Since $1-(k+1)/(\beta+k)=(\beta-1)/(\beta+k)$, the second and 
third terms in~\eqref{eq:upper-opt-error} both scale as 
$T^{-(\beta-1)/(\beta+k)}$. Moreover, $(\beta-1)/(\beta+k)<1$ 
(since $\beta>1$ and $k\geq 0$) yields 
$T^{-1}\leq T^{-(\beta-1)/(\beta+k)}$, so the initialization term 
$C_6/T$ is also dominated by $T^{-(\beta-1)/(\beta+k)}$. Combining 
and using $d\geq 1$,
\begin{equation}
\label{eq:upper-regime1}
\mathbb{E}\bigl[f(x_T)-f(x^\star)\bigr]
\leq C_9\,d\,T^{-(\beta-1)/(\beta+k)},
\end{equation}
where 
$C_9:=C_6+C_7\,\kappa^{-2(k+1)}+C_8\,\kappa^{2(\beta-1)}$.

If $1\leq T<\kappa^{2(\beta+k)}$, then $h=1$, so 
$h^{2(\beta-1)}=h^{-2(k+1)}=1$ 
and~\eqref{eq:upper-opt-error} reduces to 
$\mathbb{E}[f(x_T)-f(x^\star)]\leq(C_6+C_7\,d)/T+C_8\,d$. Since 
$T\leq\kappa^{2(\beta+k)}$, 
$T^{-(\beta-1)/(\beta+k)}\geq\kappa^{-2(\beta-1)}$. Hence, setting
\[
C_9':=\max\bigl(\kappa^{2(\beta-1)}(C_6+C_7+C_8),\;C_9\bigr),
\]
the bound~\eqref{eq:upper-regime1}, with $C_9$ replaced by $C_9'$, 
also holds in this regime. In both cases,
\begin{equation}
\label{eq:upper-final-T}
\mathbb{E}\bigl[f(x_T)-f(x^\star)\bigr]
\leq C_0\,d\,T^{-(\beta-1)/(\beta+k)},
\end{equation}
where $C_0:=C_9'$ depends only on 
$\alpha,\gamma,t_0,\kappa,L_f,G,R,C_{\mathrm{bias}},C_{\mathrm{var}}$.

\medskip
\noindent\textit{Conversion to total oracle budget $N$.}
Recall $N=2dT$, so $T=N/(2d)$. Hence
\[
T^{-(\beta-1)/(\beta+k)}
=2^{(\beta-1)/(\beta+k)}\,d^{(\beta-1)/(\beta+k)}\,
N^{-(\beta-1)/(\beta+k)},
\]
and using $1+(\beta-1)/(\beta+k)=(2\beta+k-1)/(\beta+k)$,
\[
d\,T^{-(\beta-1)/(\beta+k)}
=2^{(\beta-1)/(\beta+k)}\,d^{(2\beta+k-1)/(\beta+k)}\,
N^{-(\beta-1)/(\beta+k)}.
\]
Setting $C:=2^{(\beta-1)/(\beta+k)}\,C_0$ and combining 
with~\eqref{eq:upper-final-T},
\[
\mathbb{E}\bigl[f(x_T)-f(x^\star)\bigr]
\leq C\,d^{(2\beta+k-1)/(\beta+k)}\,N^{-(\beta-1)/(\beta+k)}.
\]
This bound is uniform over 
$g\in\mathcal{F}'_{\alpha,\beta,k}(L,L_f,G,G_g,\bar{L})$, which 
completes the proof.
\end{proof}

\end{document}
